\begin{thebibliography}{99}
\addcontentsline{toc}{chapter}{Bibliografía y fuentes consultadas}
\bibitem{GUIA2026} Universidad de León. Guía docente 2026/2027, Tecnología Aeroespacial (0710312). PDF firmado el 18 de agosto de 2026, 7 páginas. Consulta de competencias, contenidos, evaluación y bibliografía. \url{https://visor-guiadocente.unileon.es/snapshots/2026-2027/0710312}.
\bibitem{TAE0} Delgado Marcos, A.; Rubio Sierra, C.; Parra Villarejo, F. Presentación de la asignatura. ULE, 2026/2027, 4 páginas. Material docente proporcionado.
\bibitem{TAE1} Delgado Marcos, A. Historia de la aviación y carrera espacial. ULE, 2026/2027, 39 páginas. Material docente proporcionado.
\bibitem{TAE2} Delgado Marcos, A. Atmósfera estándar, atmósfera real y meteorología básica. ULE, 2026/2027, 27 páginas. Material docente proporcionado.
\bibitem{TAE3} Delgado Marcos, A. Conceptos fundamentales de la mecánica de fluidos. ULE, 2026/2027, 19 páginas. Material docente proporcionado; discrepancias documentadas.
\bibitem{TAE-P} Delgado Marcos, A. Problemas de atmósfera ISA, mecánica de fluidos y globos. ULE, 2026/2027, 6 páginas. Referencia de competencias.
\bibitem{ISA1976} NOAA, NASA y USAF. U.S. Standard Atmosphere, 1976. NASA-TM-X-74335. Consulta de constantes, capas bajas, coordenada geopotencial y tablas; no de toda la región superior. \url{https://ntrs.nasa.gov/api/citations/19770009539/downloads/19770009539.pdf}.
\bibitem{PHAK} FAA. Pilot’s Handbook of Aeronautical Knowledge, FAA-H-8083-25C, 2023. Consulta temática de construcción, vuelo, aerodinámica, sistemas, meteorología, aeropuertos y navegación. Sus referencias operativas estadounidenses no se trasladan como normativa española. \url{https://www.faa.gov/regulations_policies/handbooks_manuals/aviation/phak}.
\bibitem{FAA-AIRFRAME} FAA. Aviation Maintenance Technician Handbook: Airframe, FAA-H-8083-31B, 2023. Consulta temática de estructuras, materiales, hidráulica, electricidad, combustible, hielo, fuego y apoyo en tierra. \url{https://www.faa.gov/regulations_policies/handbooks_manuals/aviation/FAA-H-8083-31B_Aviation_Maintenance_Technician_Handbook.pdf}.
\bibitem{FAA-POWERPLANT} FAA. Aviation Maintenance Technician Handbook: Powerplant, FAA-H-8083-32B, 2023. Consulta del capítulo 1 (Aircraft Engines) y páginas oficiales de capítulos auxiliares. No se declara lectura integral de todos sus volúmenes. \url{https://www.faa.gov/regulationspolicies/handbooksmanuals/aviation/faa-h-8083-32b-aviation-maintenance-technician}.
\bibitem{NASA-BGA} NASA Glenn. Beginner’s Guide to Aeronautics. Recursos consultados sobre Wright, sustentación, resistencia inducida, compresibilidad y propulsión. URLs y alcance de las consultas en informes del proyecto. \url{https://www1.grc.nasa.gov/beginners-guide-to-aeronautics/}.
\bibitem{NASA-HISTORY} NASA. Dawn of the Space Age y Apollo 11. Consulta de hitos concretos de 1957–1958 y 1969. \url{https://www.nasa.gov/history/dawn-of-the-space-age/}.
\bibitem{NASA-ORBIT} NASA Science. Orbits and Kepler’s Laws. Consulta de leyes y descripción de órbitas. \url{https://science.nasa.gov/solar-system/orbits-and-keplers-laws/}.
\bibitem{NASA-ROCKET} NASA Glenn. Ideal Rocket Equation. Consulta de conservación de momento y ecuación ideal; el término de presión se contrasta dimensionalmente. \url{https://www1.grc.nasa.gov/beginners-guide-to-aeronautics/ideal-rocket-equation/}.
\bibitem{NASA-ISENTROP} NASA Glenn. Isentropic Flow Equations. \url{https://www.grc.nasa.gov/www/k-12/airplane/isentrop.html}. Consulta de relaciones de remanso y dominio isentrópico.
\bibitem{NASA-SHOCK} NASA Glenn. Normal Shock Wave Equations. \url{https://www.grc.nasa.gov/www/k-12/airplane/normal.html}. Consulta de saltos y pérdida de presión total.
\bibitem{NASA-WRIGHT} NASA Glenn. Wright Brothers Aircraft. \url{https://www1.grc.nasa.gov/beginners-guide-to-aeronautics/wright-brothers-aircraft/}. Consulta de ensayos y primer vuelo controlado.
\bibitem{NASA-APOLLO} NASA. Apollo 11. \url{https://www.nasa.gov/mission/apollo-11/}. Consulta de tripulación y hito lunar.
\end{thebibliography}
\chapter*{Bibliografía recomendada no consultada íntegramente}
Los títulos siguientes proceden de los metadatos de la guía oficial consultada. Verificar esa ficha no implica haber leído el texto completo. No se atribuye a un libro comercial respaldo de una ecuación cuya página no se haya examinado. La ficha pública de Alderliesten se abrió; el acceso automatizado al texto de Leishman devolvió una restricción de acceso.
\begin{itemize}
\item Pérez Sanz, L. Introducción al sistema de navegación aérea. Ibergarceta, 2013. ISBN 978-84-1545-281-2.
\item Franchini, S. y López García, Ó. Introducción a la ingeniería aeroespacial. UPM, 2008, edición citada en la guía.
\item Ledesma, M. y Baleriola, G. Meteorología aplicada a la aviación. Paraninfo, 2003. ISBN 84-283-2840-4.
\item Leishman, J. G. Introduction to Aerospace Flight Vehicles. Embry-Riddle Aeronautical University, 2023. ISBN 979-8-9852614-0-0.
\item Isidoro Carmona, A. Aerodinámica y actuaciones del avión. Paraninfo, 1999. ISBN 84-283-2640-1.
\item Muñoz, M. y Payri González, F. Motores de combustión interna alternativos. ETS Ingenieros Industriales, 1989. ISBN 84-86451-01-9.
\item Cuesta Álvarez, M. Motores de reacción. Paraninfo, 2001. ISBN 978-84-283-2825-8.
\item Moir, I. y Seabridge, A. G. Aircraft Systems. Wiley, 2009. ISBN 978-0-470-05996-8.
\item Alderliesten, R. Introduction to Aerospace Structures and Materials. TU Delft, 2018. DOI: 10.5074/t.2018.003.
\item Mouritz, A. P. Introduction to Aerospace Materials. Woodhead, 2012. ISBN 9780857095152.
\item Elices, T. Introducción a la dinámica espacial. INTA, 1991. ISBN 84-606-0382-2.
\item Anderson, J. D. Introduction to Flight. McGraw-Hill, 2005, edición citada en la guía. ISBN 0-07-123818-2.
\end{itemize}
